\documentclass[11pt,showkeys]{essay}

\begin{document}

\title{\textit{Holistic Fairness:}\\
Why Neither Equality Nor Process Is Enough}

\author{Jaxen Dutta\, \orcidlink{0009-0002-5088-4679}}
\email[]{jaxen.dutta@uwaterloo.ca}
\affiliation{David R. Cheriton School of Computer Science,\\ 
University of Waterloo, Waterloo, Ontario, Canada}

\date{January 15, 2024}

\begin{abstract}
Fairness is among the most contested yet consequential concepts in moral and political philosophy. This essay initiated its exploration by arguing whether neither egalitarian fairness, which prioritizes equal distribution of outcomes, nor procedural fairness, which emphasizes just processes, is independently sufficient to achieve justice. Eventually, it advances a synthesized account called \textit{holistic fairness}: a framework that treats equitable starting conditions and transparent, inclusive processes as mutually necessary and jointly constitutive of a just society. Drawing on philosophical, legal, and empirical literature, the essay evaluates the internal limitations of each traditional definition before defending holistic fairness as a more adaptable, culturally sensitive, and practically applicable standard for contemporary governance and social organization.
\end{abstract}

\keywords{fairness, egalitarianism, procedural justice, distributive justice, holistic fairness}

\maketitle

\section{Introduction}

Fairness is a fundamental principle deeply embedded in our societal fabric, serving as a cornerstone for ethical decision-making and just governance. At its essence, fairness refers to the impartial and equitable treatment of individuals: ensuring that all parties receive their due share of rights, opportunities, and considerations. This multifaceted concept extends its influence across domains ranging from legal systems to interpersonal relationships, from educational institutions to economic structures.

The pursuit of fairness reflects a universal aspiration for a world where justice prevails: where individuals are treated with equality and dignity irrespective of their background, identity, or circumstances. It is a guiding principle that seeks to balance power differentials, rectify injustices, and foster an environment where everyone has a fair chance to succeed.

Yet the concept resists easy definition. What counts as a fair outcome? A fair process? A fair starting point? These questions open onto deep disagreements in moral and political philosophy. This essay explores that landscape, evaluating traditional definitions before advancing a personal account that attempts to reconcile their competing strengths.

\section{Traditional Definitions}

\subsection{Egalitarian Fairness}

Among the many lenses through which fairness can be examined, the egalitarian perspective occupies a prominent place. This view holds that justice is achieved when certain goods: welfare, income, resources, wealth, and opportunity, are distributed equally or nearly equally across all members of society~\cite{temkin_equality_2019,afolayan_egalitarianism_2015}. On this account, fairness is fundamentally about the equal distribution of outcomes~\cite{eldridge_egalitarianism_2023,johnson_psychology_2023}.

Egalitarian fairness draws inspiration from the notion that equal treatment creates a level playing field, allowing each person to pursue their aspirations unencumbered by systemic inequality. Advocates argue that this approach not only promotes social harmony but also addresses historical disparities, offering a pathway to rectify past injustices~\cite{baiasu_why_2020}.

However, critics contend that absolute equality may undermine individual initiative and merit-based achievement. They argue that a more nuanced approach, one that accounts for diverse talents, contributions, and circumstances, is necessary for a truly fair society~\cite{temkin_equality_2019}. The risk is that strict distributive equality, pursued without attention to context, can itself become a form of injustice.

\subsection{Procedural Fairness}

An alternative tradition shifts the focus from outcomes to processes. Procedural fairness holds that a just society is one in which decisions are made impartially and individuals are treated with dignity, regardless of the final results~\cite{dahlstrom_what_nd,mushtaq_basic_2020}. This perspective underscores the importance of transparent, accountable systems in which rules are applied consistently and individuals have a meaningful voice in decision-making~\cite{canada_what_2017}.

The emphasis falls on the integrity of the mechanisms governing social interaction rather than on the outcomes those mechanisms produce. When processes are fair, proceduralists argue, people are more likely to accept outcomes, even unfavourable ones, as legitimate.

Yet critics observe that procedural fairness alone cannot address underlying structural inequalities. Fair processes applied to an unequal starting configuration will reproduce and entrench existing hierarchies~\cite{macdonald_judicial_1980,galligan_procedural_1997}. A perfectly transparent process can yield systematically biased outcomes when underlying conditions are unequal, pointing toward the need for a framework that integrates both procedural and distributive considerations~\cite{ferguson_procedural_2014,rescher_review_2003,morse_ends_2022}.

\section{Holistic Fairness: A Personal Definition}

Reflecting on these two traditions and their respective limitations, I find myself most aligned with what I call \textit{holistic fairness}: a synthesized framework that treats both equal opportunities and just processes as necessary components of a fair society. Holistic fairness recognizes that genuine justice requires equitable starting conditions \textit{and} respectful, inclusive mechanisms for adjudicating outcomes. Neither is sufficient on its own.

\subsection{Alignment with Personal Values}

Holistic fairness aligns with my core conviction that every individual possesses intrinsic worth that transcends their social position or accumulated resources. It moves beyond crude redistribution to ensure that each person has a genuine opportunity to participate in and contribute to society. At the same time, it honours the proceduralist insight that how decisions are made matters morally, and that transparency, accountability, and participatory governance are not mere formalities but constitutive elements of a just order.

\subsection{Practical Application}

The practical purchase of holistic fairness is one of its principal virtues. The framework applies coherently across a range of institutional contexts: from advocating for equal access to quality education, to designing transparent and inclusive workplace hiring practices, to reforming criminal justice systems, to developing social welfare policies that address structural disadvantage, besides just the ones that affect individual cases. Its adaptability is deliberate: holistic fairness does not prescribe a single mechanism for achieving justice in every context, but provides evaluative criteria against which particular arrangements can be assessed and refined.

\subsection{Cultural and Societal Influences}

My alignment with holistic fairness has been shaped by exposure to diverse cultural contexts, a historical awareness of how formal equality has often masked substantive inequality, and engagement with social movements that have insisted on the inseparability of procedural and distributive justice. Holistic fairness, in this sense, is not only a philosophical position but a response to lived experience and collective struggle.

\section{Comparison with Alternatives}

\subsection{Limitations of Established Definitions}

Returning to the traditional frameworks with holistic fairness as a reference point clarifies their respective weaknesses. Egalitarian fairness, by focusing exclusively on distributional outcomes, risks homogenizing conditions in ways that suppress individual agency and fail to account for the fact that equal treatment of unequal starting positions perpetuates inequality~\cite{temkin_equality_2019,afolayan_egalitarianism_2015}. Procedural fairness, by privileging the integrity of processes over their results, can become complicit in the reproduction of structural injustice: a perfectly transparent process can yield systematically biased outcomes when underlying conditions are unequal~\cite{macdonald_judicial_1980,galligan_procedural_1997}.

Neither framework, taken in isolation, provides the normative resources needed to diagnose and remedy the full spectrum of injustices present in contemporary societies.

\subsection{Advantages of Holistic Fairness}

Holistic fairness overcomes these limitations by insisting on the mutual dependence of distributional and procedural considerations. Its key advantages are threefold. First, it is \textit{comprehensive}: by attending to both starting conditions and decision-making processes, it captures a wider range of injustices than either alternative alone. Second, it is \textit{adaptive}: the framework provides evaluative criteria applicable flexibly across diverse contexts, cultures, and historical moments. Third, it is \textit{integrative}: rather than forcing a choice between individual initiative and collective well-being, holistic fairness recognizes these as complementary rather than competing goods.

\section{Conclusion}

Fairness is neither a fixed standard nor a simple principle. It is a dynamic, contested, and contextually sensitive idea whose meaning is continuously renegotiated in light of evolving social arrangements, cultural norms, and moral understandings. The diversity of philosophical definitions: egalitarian, procedural, and beyond, reflects the genuine complexity of the concept rather than a failure to arrive at the right answer.

Within this landscape, holistic fairness represents my considered position: a synthesis that refuses to treat distributional equality and procedural integrity as mutually exclusive, insisting instead on their joint necessity. The framework is adaptable, culturally sensitive, and practically applicable. These are qualities essential for any account of fairness capable of guiding real institutions in a diverse and rapidly changing world.

Ultimately, the subjectivity of fairness is not a weakness to be overcome but a feature to be embraced. By acknowledging that our conceptions of justice are shaped by personal values, cultural inheritances, and social experiences, we create space for the kind of genuine dialogue and mutual understanding that justice itself requires.

\bibliography{holistic-fairness}
\end{document}
